\documentclass[11pt]{article}

\usepackage[a4paper,margin=1.05in]{geometry}
\usepackage{mathpazo}
\usepackage[T1]{fontenc}
\usepackage{amsmath,amssymb}
\usepackage{microtype}
\usepackage{booktabs}
\usepackage{tabularx}
\usepackage{enumitem}
\usepackage{parskip}
\usepackage{titlesec}
\usepackage{fancyhdr}
\usepackage{tcolorbox}
\usepackage[colorlinks=true,linkcolor=accent,urlcolor=accent,citecolor=accent]{hyperref}

\tcbuselibrary{skins,breakable}

\definecolor{accent}{HTML}{1F5C6B}
\definecolor{accentbg}{HTML}{DCE9EC}
\definecolor{softgray}{HTML}{6B6B6B}
\definecolor{warnrule}{HTML}{B8860B}
\definecolor{warnbg}{HTML}{FBF0D9}

\titleformat{\section}{\large\bfseries\color{accent}}{\thesection}{0.7em}{}
\setlist[itemize]{leftmargin=1.3em,itemsep=2pt,topsep=3pt}

\pagestyle{fancy}
\fancyhf{}
\renewcommand{\headrulewidth}{0.3pt}
\lhead{\small\color{softgray}The model: a draft}
\rhead{\small\color{softgray}\thepage}

\newtcolorbox{notebox}{
  enhanced, breakable, colback=accentbg!55, colframe=accent,
  boxrule=0pt, leftrule=2.5pt, arc=1pt, left=10pt, right=10pt, top=8pt, bottom=8pt}

\newcommand{\Ncal}{\mathcal{N}}
\newcommand{\Mset}{\mathcal{M}}
\newcommand{\Dset}{\mathcal{D}}
\newcommand{\Rset}{\mathcal{R}}
\newcommand{\Tr}{^{\!\top}}
\DeclareMathOperator{\diag}{diag}
\DeclareMathOperator*{\Cov}{Cov}
\DeclareMathOperator{\Corr}{corr}

\title{\bfseries The model: a draft}
\author{}
\date{}

\begin{document}
\maketitle
\vspace{-2.5em}

\section{Indices and dimensions}

\begin{center}\small
\begin{tabularx}{\textwidth}{@{}llX@{}}
\toprule
 & Range & \\
\midrule
$j$ & $1,\dots,P$ & model parameters \\
$q$ & $1,\dots,Q$ & model species \\
$r$ & $1,\dots,M$ & readouts \\
$s$ & $1,\dots,S$ & scenarios \\
$c$ & $1,\dots,C$ & study cohorts; $s(c)$ is the scenario cohort $c$ was run under \\
$k$ & $1,\dots,K_{rc}$ & statistics that cohort $c$ reports for readout $r$ \\
$i$ & $1,\dots,N$ & simulated patients \\
$\Rset_c$ & $\subseteq\{1,\dots,M\}$ & the readouts that cohort $c$ reports \\
$K_c$ & $\sum_{r\in\Rset_c}K_{rc}$ & the rows that cohort $c$ contributes \\
\bottomrule
\end{tabularx}
\end{center}

\noindent
$\theta\in\mathbb{R}^P_{+}$ are the parameters and $\vartheta=\log\theta$. A
\emph{target} is a $(r,c)$ pair. Readouts are on the log scale. $n_c$ is the real
cohort size, as from the actual data it comes from. A target names a row, not a
unit of independence; the cohort is that.

\section{The model}

\paragraph{Population.}
\begin{equation}
\vartheta_i \mid \mu,\omega \;\sim\; \Ncal_P\big(\mu,\ \Omega\big),
\qquad
\Omega \;=\; D_\omega\,R\,D_\omega,
\qquad
D_\omega=\diag(\omega),
\label{eq:pop}
\end{equation}
with $R$ a fixed correlation matrix. Realised with common random numbers, so
that $\vartheta_i$ is a smooth function of $(\mu,\omega)$:
\begin{equation}
\vartheta_i \;=\; \mu + D_\omega L_R z_i,
\qquad
z_i \stackrel{\text{iid}}{\sim} \Ncal_P(0,I_P) \ \text{ drawn once},
\qquad
R = L_R L_R\Tr .
\label{eq:crn}
\end{equation}

\paragraph{Mechanism.} Simulate patient $i$ under scenario $s$ and read off the
species at the measured time:
\begin{equation}
y_i(s) \;=\; g\big(\vartheta_i;\,s\big) \;\in\; \mathbb{R}^Q_{+} .
\label{eq:mech}
\end{equation}
A readout composes those species into the scalar a study reports, on the log
scale:
\begin{equation}
x_{irc} \;=\; h_r\big(y_i(s(c))\big).
\label{eq:readout}
\end{equation}

\paragraph{Discrepancy.} Two kinds of error, distinguished by where they enter
and therefore by how they propagate:
\begin{equation}
\tilde x_{irc}
\;=\;
\kappa_r\;h_r\big(
\underbrace{e^{\beta}\odot y_i(s(c))}_{\text{mechanism, upstream of }h_r}
\big)
\;+\;
\gamma_r ,
\qquad
\underbrace{\gamma_r = Z_r\Tr a,
\quad
\log\kappa_r = Z_r\Tr b}_{\text{measurement, downstream of }h_r} .
\label{eq:disc}
\end{equation}
$\gamma$ and $\kappa$ are the location and scale halves of a linear map on the
readout.

$\beta\in\mathbb{R}^Q$ is a log-scale bias on each \emph{species}, shared across
readouts and scenarios. It enters upstream of $h_r$, so it propagates through the
composition on its own: a ratio inherits $\beta_{q_1}-\beta_{q_2}$, a percentage
inherits $\beta_q - \beta_{\text{nuc}}$, a sum inherits a weighted combination.

$a$ and $b$ are measurement effects on the \emph{readout}, entering downstream of
$h_r$, so they do not propagate. $Z_r$ holds attributes of the readout that have
no mechanistic counterpart: assay modality, and whether the quantity is a
density, a fraction or a concentration.

$\beta=0$, $a=0$, $b=0$ is a model with no discrepancy, that is, a model that is not biased in scale or location on any of the observed species.

\paragraph{The cohort vector of one patient.} Stack the corrected readouts that
cohort $c$ reports, for a single patient:
\begin{equation}
\tilde x_{ic} \;=\; \big(\tilde x_{irc}\big)_{r\in\Rset_c}
\;\in\; \mathbb{R}^{|\Rset_c|} .
\label{eq:stack}
\end{equation}
Built from $\tilde x$ and not $x$: $\beta$ enters inside $h_r$, so
\eqref{eq:disc} does not factor through \eqref{eq:readout}.

\paragraph{Cohort selection.} Where cohort $c$ reports an eligibility criterion,
$w^{(c)}_i\in[0,1]$ is a fixed smooth indicator of whether patient $i$ meets it,
and $w^{(c)}_i \equiv 1$ otherwise. Write $\tilde F_{c}$ for the
$w^{(c)}$-weighted empirical distribution of $\{\tilde x_{ic}\}_{i=1}^N$ on
$\mathbb{R}^{|\Rset_c|}$, and $\tilde F_{rc}$ for its $r$-th marginal.

\paragraph{Statistics.} Cohort $c$ reports statistics $T_{rc1},\dots,T_{rcK_{rc}}$
for readout $r$, each a functional of a sample. The model computes the same
functionals of its own predicted sample:
\begin{equation}
\tau_{rck}(\phi) \;=\; T_{rck}\big(\tilde F_{rc}\big).
\label{eq:stat}
\end{equation}
Order the rows of cohort $c$ by $r\in\Rset_c$ and then by $k$ within $r$, and
stack them into $\tau_c(\phi)\in\mathbb{R}^{K_c}$, and the reported values the
same way into $\widehat T_c$. Since $\tilde F_{rc}$ is a marginal of $\tilde
F_c$, the joint form changes $V_c$ and leaves $\tau$ alone.

\paragraph{Observation.} One cohort is one block:
\begin{equation}
\widehat T_{c} \;\sim\; \Ncal_{K_{c}}\big(\tau_{c}(\phi),\ V_{c}\big),
\qquad
\widehat T_{c} \perp \widehat T_{c'} \ \text{ for } c\neq c' .
\label{eq:obs}
\end{equation}
Under the ordering of \eqref{eq:stat}, the diagonal blocks of $V_c$ are the
per-target covariances and the off-diagonal blocks carry the dependence between
readouts of one cohort.

\paragraph{Sampling covariance.} $V_{c}$ is the covariance of the reported
statistics of cohort $c$ under resampling of its $n_c$ patients. It has two
ingredients. The first is a bootstrap from the predicted population at a plug-in
$\phi_0$, held fixed thereafter:
\begin{equation}
V^{\text{boot}}_{c}
\;=\;
\Cov_{\ast}\!\Big[\,T_{c}\big(\tilde F^{\ast}_{c}\big)\,\Big]_{\phi_0},
\qquad
\tilde F^{\ast}_{c}
= \text{resample of } n_c \text{ patients from } \tilde F_{c},
\label{eq:V}
\end{equation}
where a resampled patient carries its whole vector $\tilde x_{ic}$ from
\eqref{eq:stack}. Drawing patients rather than readout values is what moves the
cohort's dependence structure into $V_c$.

The second is $E_c$, from the emulator $\widehat g$ standing in for $g$. Define
it on the statistics and not on the readouts, because the map from a readout
error to a statistic error depends on the functional. Draw $L$ held-out clouds
$z^{(\ell)}$ at $\phi_0$, evaluate each both ways, and take the covariance of the
difference:
\begin{equation}
d^{(\ell)}_c
\;=\;
\tau^{\widehat g}_c\big(\phi_0;\,z^{(\ell)}\big)
\;-\;
\tau^{g}_c\big(\phi_0;\,z^{(\ell)}\big),
\qquad
E_c \;=\; \Cov_{\ell}\big[\,d^{(\ell)}_c\,\big],
\qquad \ell=1,\dots,L .
\label{eq:Ec}
\end{equation}
This is $K_c\times K_c$ by construction. Report $\bar d_c$ beside it: a constant
part of the emulator error is a fixed offset that $\gamma_r$ absorbs, so only
the varying part belongs in $V_c$.

The fallback is $E_c=\diag(\varepsilon^2)$, with $\varepsilon_r$ the held-out
residual of readout $r$ repeated across its $K_{rc}$ statistics. It asserts that
the emulator errs independently across readouts, which readouts sharing a species
do not satisfy, and that a readout error enters a scale statistic with the same
loading as a location statistic, where the true loading is near zero.

\paragraph{Assembling $V_c$.} A reported uncertainty measures $\sqrt{V_c}$ on its
own diagonal, from the real patients, and carries no dependence on $\phi_0$.
Split the bootstrap into correlations and scales, and replace only the scales:
\begin{equation}
V_c \;=\; D_c\,\rho_c\,D_c \;+\; E_c,
\qquad
\rho_c=\Corr\big(V^{\text{boot}}_c\big),
\qquad
D_c=\diag(d_c),
\qquad
d_{c,rck}=
\begin{cases}
U_{rck}, & \text{reported,}\\[2pt]
\sqrt{\big(V^{\text{boot}}_c\big)_{rck,rck}}, & \text{otherwise,}
\end{cases}
\label{eq:Vsplit}
\end{equation}
with $U_{rck}$ the reported uncertainty of row $(r,c,k)$, on the log scale. No
source reports a correlation between readouts, so $\rho_c$ always comes from the
cloud. This is the split \eqref{eq:pop} already makes for the population
covariance.

The ratio
\begin{equation}
\varrho_{rck} \;=\; U_{rck}\Big/\sqrt{\big(V^{\text{boot}}_c\big)_{rck,rck}}
\label{eq:ratio}
\end{equation}
is computable at $\phi_0$ before anything is fitted. Above one, the predicted
cloud is too narrow at that readout.

\paragraph{One number, one job.} A standard deviation is a population spread and
belongs with the scale statistics of \eqref{eq:stat}. A standard error is a
sampling uncertainty and belongs in $D_c$. At $n_c=10$ they differ by a factor of
$3.2$, so the extraction has to record which one the source printed.

A source reporting only an uncertainty can supply either, not both. Send it to
whichever is scarcer on that row. Where a scale statistic already exists, $U$
goes to $D_c$. Where none does, $U$ is the only statement about that width, so
make it a scale statistic: the sampling standard deviation of a median at size
$n_c$ is a functional of $\tilde F_c$, so \eqref{eq:stat} predicts it by the
bootstrap of \eqref{eq:V} evaluated at $\phi$ rather than frozen at $\phi_0$, and
\eqref{eq:ratio} becomes a residual. Freeze the resample indices alongside
$z_{1:N}$ so the row stays smooth in $\phi$.

Such a row carries roughly $1/\sqrt{2(n_c-1)}$ on the log scale, the precision
\eqref{eq:omegameas} gives a measured spread, and it is approximate because it
depends on how the source built its interval. It buys precision on $\omega$
without breaking the alias: $\kappa_r$ scales a sampling standard deviation
exactly as it scales an interquartile range.

We may need to try to evaluate this covariance at other $\phi_0$, in order to ascertain if a plug-in estimator is ok to use.
Section \ref{sec:flat} supplies the plug-in itself, at
$\phi_0=(\hat\mu_{\text{flat}},\ \omega_0)$.

\section{Priors}

\paragraph{Location.}
\begin{equation}
\mu \;\sim\; \Ncal_P\big(\mu_0,\ \Sigma_1\big),
\label{eq:mu}
\end{equation}
$\mu_0$ and $\Sigma_1$ the mean and covariance of the stage-1 composite prior on
the log scale. Recall that this is the stage 1- fitting in which we are specifically trying to learn the location parameter per QSP parameter.

\paragraph{Spread.} Let $\Mset$ be the parameters whose between-patient spread a
source has measured, with sample size $n_j$. For $j\in\Mset$,
\begin{equation}
\log\omega_j \;\sim\; \Ncal\!\left(\log\omega_{0j},\ \frac{1}{2(n_j-1)+\tau^{-2}}\right).
\label{eq:omegameas}
\end{equation}
For $j\notin\Mset$,
\begin{equation}
\log\omega_j \;=\; \log\omega_{0j} + s + u_j,
\qquad
s\sim\Ncal(0,\tau_s^2),
\qquad
u_j\sim\Ncal(0,\tau_\omega^2),
\qquad
\sum_{j\notin\Mset} u_j = 0 .
\label{eq:omegaassumed}
\end{equation}

\paragraph{Measurement discrepancy.} Free, because it has no parameter twin.
$Z$ has an intercept, so $a_1$ and $b_1$ are the global location and width
corrections.
\begin{equation}
a \;\sim\; \Ncal\big(0,\ \sigma_a^2 I\big),
\qquad
b \;\sim\; \Ncal\big(0,\ \sigma_b^2 I\big),
\qquad
\sigma_a=\sigma_b=0.5 .
\label{eq:abprior}
\end{equation}

\paragraph{Mechanism discrepancy.} Tightly shrunk, because it competes with
$\mu$:
\begin{equation}
\beta_q \;=\; \tau_\beta\,\beta^{\text{raw}}_q,
\qquad
\beta^{\text{raw}}_q \sim \Ncal(0,1),
\qquad
\log\tau_\beta \sim \Ncal\big(\log 0.15,\ 0.5^2\big).
\label{eq:betaprior}
\end{equation}
$\beta_q$ is set to zero for every species outside a declared subset
$\mathcal{S}\subset\{1,\dots,Q\}$. The default $\mathcal{S}$ is the shared
denominators, $\mathrm{nucleated\_total}$ and $V_T$.

\section{Posterior and deliverable}

\begin{equation}
\phi \;=\;
\big(\mu,\ s,\ u^{\text{raw}},\ \log\omega_{\Mset},\
a,\ b,\ \beta^{\text{raw}}_{\mathcal{S}},\ \log\tau_\beta\big).
\end{equation}

\begin{equation}
p(\phi\mid\Dset)
\;\propto\;
p(\phi)\;\prod_{c=1}^{C}\;
\Ncal\big(\widehat T_{c};\ \tau_{c}(\phi),\ V_{c}\big).
\label{eq:post}
\end{equation}

The virtual population is the posterior predictive over patients,
\begin{equation}
p(\theta_{\text{new}}\mid\Dset)
\;=\;
\int \Ncal_P\!\big(\log\theta_{\text{new}};\ \mu,\ \Omega\big)\;
p(\phi\mid\Dset)\;d\phi ,
\label{eq:vpop}
\end{equation}
and it is emitted as an ensemble: one population per posterior draw of $\phi$,
so that between-patient variability and uncertainty about $\phi$ stay separable
in whatever is computed downstream.

\section{The flat fit}
\label{sec:flat}

The flat fit estimates the centre of the population and holds its shape fixed.
Set $\omega=\omega_0$ and drop the three spread terms from the parameter vector:
\begin{equation}
\phi_{\text{flat}} \;=\;
\big(\mu,\ a,\ b,\ \beta^{\text{raw}}_{\mathcal{S}},\ \log\tau_\beta\big).
\label{eq:phiflat}
\end{equation}
Nothing else changes. Equations \eqref{eq:crn} to \eqref{eq:post} run as
written, with $\omega$ a constant rather than a function of $\phi$.

\paragraph{Why $\omega_0$ and not zero.} Setting $\omega=0$ collapses the cloud
onto one point, and three things die with it. $\tilde F_c$ becomes a point mass,
so \eqref{eq:V} returns zero and $V_c$ has to come from outside the model.
$w^{(c)}$ becomes a switch on the whole cohort instead of on patients, so the
eligibility rule stops working and turns discontinuous in $\mu$. And every
location functional of a point mass returns the same number, so a cohort that
reports a mean beside a median predicts one value for two rows.

Each of those is a diagnostic, and diagnostics are what the flat fit is for. The
third is the sharpest: the gap between a reported mean and a reported median is
the skewness of the pushforward, which at $\omega=\omega_0$ the model predicts
rather than absorbs.

\paragraph{What it buys.}

\begin{center}\small
\begin{tabularx}{\textwidth}{@{}lX@{}}
\toprule
$V_c$ is computable & \eqref{eq:V} resamples a real cloud, so the bootstrap works unchanged. \\
The scale rows are held out & Predict them at $\omega_0$ and report the ratio per readout. That is the width shortfall, before any population parameter is estimated. \\
$\hat\mu$ needs no translation & It is already a population centre, so there is no Jensen gap to correct afterwards. \\
It supplies the plug-in & Take $\phi_0=(\hat\mu_{\text{flat}},\ \omega_0)$ for \eqref{eq:V}. \\
\bottomrule
\end{tabularx}
\end{center}

Fit the location rows and hold the scale rows out, so the ratio above is an
honest comparison. Refitting with them included is the natural second run, and it
is what identifies $b$.

\paragraph{Cost.} A gradient step still costs $C\times N$ forward passes, but $N$
only has to resolve the location statistics and a rough width, not the flanks at
a precision that supports inference on $\omega$, so it sits far below the $N$ of
the population fit.

\paragraph{What it cannot test.} The shape of the population: $\omega$, $R$, and
the direction $s$ and $b_1$ share. It does test the prior on $\mu$, the forward
map, the discrepancy terms, the block structure of \eqref{eq:obs} and the
sampler, and a fault in any of those is one the population fit would also have.

\section{Fixed inputs}

\begin{center}\small
\begin{tabularx}{\textwidth}{@{}lX@{}}
\toprule
$R$ & Population correlation matrix. Not identified by these data. \\
$\omega_0$ & Prior centre of the spreads, per parameter. \\
$\mu_0,\ \Sigma_1$ & Stage-1 composite prior, log scale. \\
$h_r$ & Readout composition, including unit conversions. Already in the target definitions. \\
$Z$ & Readout attributes for the measurement discrepancy: assay modality, quantity type. \\
$\mathcal{S}$ & Species allowed a mechanism discrepancy. Default: the shared denominators. \\
$z_{1:N}$ & The latent cloud, drawn once. \\
$w^{(c)}$ & Eligibility weights, read from each paper's criterion. \\
$V_{c}$ & Assembled by \eqref{eq:Vsplit} at $\phi_0$, then frozen. One block per cohort. \\
$E_c$ & Emulator error covariance on the statistics of cohort $c$, from \eqref{eq:Ec}. \\
$U_{rck}$ & Reported uncertainty of a statistic, where a source gives one. Log scale. \\
$\tau_s,\ \tau_\omega,\ \tau,\ \sigma_a,\ \sigma_b$ & Prior widths. \\
\bottomrule
\end{tabularx}
\end{center}

\section{Known simplifications}

\begin{itemize}
\item \textbf{Independence across cohorts.} Equation \eqref{eq:obs} is block
diagonal by cohort, so what remains asserted is that distinct cohorts are
distinct patients. That holds across studies. It fails where two cohorts come
from one publication and share people, as when a subgroup is reported beside the
whole. Check it before splitting any pooled target: a split that manufactures two
cohorts from one set of patients moves the problem instead of removing it.

Report the effective row count $K_c/\{1+(K_c-1)\bar\rho_c\}$ at mean off-diagonal
correlation $\bar\rho_c$. The source supplying 19 of the 49 targets, from about
ten patients per arm, is the one to look at first.

\item \textbf{No study offset.} A free $\eta_c$ would be aliased with the
measurement discrepancy. It is omitted and $\gamma_r$ absorbs the study effect,
so $\gamma_r$ mixes assay bias with study bias.

\item \textbf{$V$ frozen at $\phi_0$.} Equation \eqref{eq:post} is an estimating
equation, not a likelihood. Letting $V$ move with $\phi$ would reintroduce a
$-\tfrac12\log\det V$ term that rewards a narrow population regardless of fit.

\item \textbf{$s$ and $b_1$ share a direction.} To first order every scale
statistic depends on the global width multiplier and the global width correction
only through their sum, so the split between a narrow simulator and a narrow
population is decided by $\tau_s$ and $\sigma_b$ unless some readout is driven
mainly by $\Mset$. Report the joint posterior of the pair.

\item \textbf{$Z$ is asserted, not tested.} The claim that assay modality and
quantity type have no mechanistic twin, while species-level bias does, is an
argument rather than a result. The test is to project each column of $Z$ onto
$\mathrm{col}\big(\partial\tau/\partial(\mu,\omega)\big)$ and read the residual
norm. That needs the emulator, and it has not been run.
\end{itemize}

\end{document}
